Il peut arriver que certains capteurs, comme le GPS ou l’IMU, donnent des valeurs étranges à cause d’un signal instable, d’un environnement difficile (par exemple: immeubles, tunnel), ou même de bugs ponctuels. Pour éviter que ces erreurs ne faussent les résultats, nous avons mis en place un petit script qui détecte et enlève automatiquement les valeurs incohérentes.

Le fichier \texttt{valeurs\_aberrantes.py} utilise une méthode statistique basée sur l’IQR (intervalle interquartile). En effet, il regarde l’altitude sur une petite fenêtre autour de chaque point, et si une valeur est trop différente de ses voisines, elle est considérée comme aberrante et supprimée. C’est utile pour supprimer des pics d’altitude qui n’ont pas de sens.

Ce nettoyage se fait automatiquement avant l’analyse des performances, donc nous n’avons pas besoin de le faire à la main. Le fichier nettoyé est ensuite utilisé normalement avec le reste du programme.

Grâce à ça, même si un capteur dysfonctionne un peu à un moment, on peut tout de même analyser le trajet proprement, sans que tout soit faussé. C’est une sécurité en plus qui rend le système plus fiable, surtout quand les conditions ne sont pas idéales.

